\chapter{Representations of Groups and Lie Algebras}\label{representation-theory}
\chapterstatus{complete}

\section{Introduction}\label{representation-theory:section-introduction}

A representation of a group or a Lie algebra is an action by linear maps. This chapter develops the
basic theory: Schur's lemma, Maschke's theorem and character theory for finite groups, nilpotent and solvable Lie algebras with the
theorems of Engel and Lie, and the complete
classification of the finite-dimensional representations of the Lie algebra $\mathfrak{sl}_2(\CC)$. The last
result is used repeatedly in Hodge theory: on the cohomology of a compact K\"ahler manifold the
Lefschetz operator, its adjoint and the degree operator span a copy of $\mathfrak{sl}_2$, and the hard
Lefschetz theorem and the primitive decomposition are consequences of the structure of
$\mathfrak{sl}_2$-representations (Chapters~\ref{kahler} and~\ref{lefschetz-theorems}).

References are \cite{serre-linear-representations}, \cite{fulton-harris} and \cite{humphreys-lie}.

\noindent\textbf{Prerequisites.} Chapter~\ref{groups} (Groups), Chapter~\ref{fields} (Fields and Galois
Theory), Chapter~\ref{linear-algebra} (Linear Algebra) and Chapter~\ref{multilinear-algebra}
(Multilinear Algebra).

Throughout, $k$ is a field and representations are on $k$-vector spaces.

\section{Representations of groups}\label{representation-theory:section-group-representations}

\begin{definition}\label{representation-theory:definition-representation}
A \emph{representation} of a group $G$ is a vector space $V$ with a homomorphism
$\rho \colon G \to \GL(V)$; we write $gv$ for $\rho(g)v$. A \emph{homomorphism} of representations (or
$G$-equivariant map) is a linear map $f$ with $f(gv) = gf(v)$; the space of such maps is $\Hom_G(V, W)$. A
\emph{subrepresentation} is a $G$-stable subspace. $V$ is \emph{irreducible} if $V \neq 0$ and its only
subrepresentations are $0$ and $V$, and \emph{completely reducible} (or \emph{semisimple}) if it is a direct
sum of irreducible subrepresentations. The \emph{invariants} are $V^G = \{v : gv = v \text{ for all } g\}$.
\end{definition}

\begin{example}\label{representation-theory:example-constructions}
Given representations $V$, $W$: quotients $V/U$ by subrepresentations, $V \oplus W$, $V \otimes W$ with
$g(v \otimes w) = gv \otimes gw$, the dual $V^*$ with $(g\varphi)(v) = \varphi(g^{-1}v)$, and $\Hom_k(V, W)$ with
$(gf)(v) = g f(g^{-1}v)$ are representations, and $\Hom_k(V, W)^G = \Hom_G(V, W)$. The \emph{trivial
representation} is $k$ with $gv = v$. If $G$ acts on a finite set $X$, the \emph{permutation representation} is
$k^X$ with $g e_x = e_{gx}$; for $X = G$ with the left action this is the \emph{regular representation}.
\end{example}

\begin{lemma}[Schur]\label{representation-theory:lemma-schur}
Let $V$, $W$ be irreducible representations. Every $f \in \Hom_G(V, W)$ is zero or an isomorphism. If $k$ is
algebraically closed and $V$ is finite-dimensional, then $\Hom_G(V, V) = k \cdot \id$.
\end{lemma}

\begin{proof}
$\ker f$ and $\im f$ are subrepresentations, so $\ker f \in \{0, V\}$ and $\im f \in \{0, W\}$. If $k$ is algebraically
closed and $f \in \Hom_G(V, V)$, $f$ has an eigenvalue $\lambda$, and $f - \lambda \in \Hom_G(V, V)$ is not an isomorphism,
hence zero.
\end{proof}

\begin{theorem}[Maschke]\label{representation-theory:theorem-maschke}
Let $G$ be finite with $\operatorname{char} k \nmid |G|$. Every subrepresentation $U$ of a representation $V$ has a
$G$-stable complement. Hence every finite-dimensional representation is completely reducible.
\end{theorem}

\begin{proof}
Choose a linear projection $\pi \colon V \to U$ (a linear map with $\pi|_U = \id$), using a complement of $U$
(Proposition~\ref{linear-algebra:proposition-subspaces}), and put
$\pi_0 = \frac{1}{|G|}\sum_{g \in G} g\pi g^{-1}$. It is $G$-equivariant, has image in $U$, and restricts to the identity
on $U$; so $\ker\pi_0$ is a $G$-stable complement of $U$. The second statement follows by induction on the
dimension.
\end{proof}

\begin{counterexample}\label{representation-theory:counterexample-not-semisimple}
Both hypotheses of Maschke's theorem are needed. The group $\ZZ$ acts on $\CC^2$ by
$n \mapsto \begin{pmatrix} 1 & n \\ 0 & 1 \end{pmatrix}$; the line $\CC e_1$ is stable, but it has no stable complement, since a
stable line would be spanned by an eigenvector and $e_1$ spans the only eigenline. Similarly $\ZZ/p\ZZ$ acts on
$\mathbb{F}_p^2$ by the same matrices (with $n$ read modulo $p$), and the representation is not completely
reducible although $\ZZ/p\ZZ$ is finite; here $p$ divides the order of the group.
\end{counterexample}

\section{Characters of finite groups}\label{representation-theory:section-characters}

In this section $G$ is finite, $k = \CC$, and representations are finite-dimensional.

\begin{definition}\label{representation-theory:definition-character}
The \emph{character} of a representation $V$ is $\chi_V \colon G \to \CC$, $\chi_V(g) = \operatorname{tr}(g|_V)$. A
\emph{class function} is a function $G \to \CC$ constant on conjugacy classes. On class functions put the
Hermitian inner product $\langle \phi, \psi \rangle = \frac{1}{|G|}\sum_{g \in G}\phi(g)\overline{\psi(g)}$.
\end{definition}

\begin{proposition}\label{representation-theory:proposition-character-properties}
Characters are class functions, $\chi_V(e) = \dim V$, $\chi_{V \oplus W} = \chi_V + \chi_W$,
$\chi_{V \otimes W} = \chi_V\chi_W$, and $\chi_{V^*}(g) = \overline{\chi_V(g)} = \chi_V(g^{-1})$. Moreover
$\dim V^G = \frac{1}{|G|}\sum_g \chi_V(g)$.
\end{proposition}

\begin{proof}
The trace is invariant under conjugation, additive on direct sums and multiplicative on tensor products
(compute in bases). Each $g$ has finite order, so $g|_V$ satisfies $x^n - 1$ and is diagonalisable with roots of
unity $\lambda_i$ as eigenvalues (Theorem~\ref{linear-algebra:theorem-diagonalisable}); then
$\chi_{V^*}(g) = \sum \lambda_i^{-1} = \sum \bar\lambda_i$. Finally $p = \frac{1}{|G|}\sum_g g$ is a projection onto $V^G$
(it is the identity on $V^G$ and $hp = p$ for all $h$), and the trace of a projection is the dimension of its
image.
\end{proof}

\begin{theorem}[Orthogonality]\label{representation-theory:theorem-orthogonality}
For representations $V$, $W$, $\langle \chi_W, \chi_V \rangle = \dim \Hom_G(V, W)$. Consequently the characters of
irreducible representations are orthonormal, a representation is determined up to isomorphism by its
character, and $V$ is irreducible if and only if $\langle \chi_V, \chi_V \rangle = 1$.
\end{theorem}

\begin{proof}
$\Hom_G(V, W) = \Hom_\CC(V, W)^G$, and $\Hom_\CC(V, W) \cong V^* \otimes W$ as representations, so by
Proposition~\ref{representation-theory:proposition-character-properties} its dimension is
$\frac{1}{|G|}\sum_g \overline{\chi_V(g)}\chi_W(g)$. For irreducible $V$, $W$, Schur's lemma gives
$\dim\Hom_G(V, W) = 1$ if $V \cong W$ and $0$ otherwise. By Maschke's theorem $V \cong \bigoplus_i V_i^{m_i}$ with
distinct irreducible $V_i$, and $m_i = \langle \chi_V, \chi_{V_i} \rangle$ is determined by $\chi_V$; also
$\langle \chi_V, \chi_V \rangle = \sum m_i^2$.
\end{proof}

\begin{theorem}\label{representation-theory:theorem-number-irreducibles}
The characters of the irreducible representations of $G$ form an orthonormal basis of the space of class
functions. In particular the number of isomorphism classes of irreducible representations equals the number
of conjugacy classes, and if $V_1, \ldots, V_r$ are the irreducible representations, then
$\sum_i (\dim V_i)^2 = |G|$.
\end{theorem}

\begin{proof}
The regular representation has character $|G|$ at $e$ and $0$ elsewhere, so by
Theorem~\ref{representation-theory:theorem-orthogonality} it contains each $V_i$ with multiplicity
$\langle \chi_{\mathrm{reg}}, \chi_{V_i}\rangle = \dim V_i$; comparing dimensions gives $\sum(\dim V_i)^2 = |G|$. In particular there
are finitely many irreducible representations up to isomorphism.

It remains to show that a class function $\phi$ orthogonal to all $\chi_{V_i}$ is zero. For a representation $V$ let
$T_V = \sum_g \overline{\phi(g)}\, g|_V$. Since $\phi$ is a class function, $T_V$ commutes with all $h|_V$. If $V$ is
irreducible, $T_V = \lambda \id$ by Schur's lemma, and taking traces,
$\lambda \dim V = \sum_g \overline{\phi(g)}\chi_V(g) = |G|\langle \chi_V, \phi \rangle = 0$. So $T_V = 0$ for every irreducible $V$, hence
for every $V$ by Maschke's theorem, in particular for the regular representation; but there
$T(e_e) = \sum_g \overline{\phi(g)} e_g$, so $\phi = 0$.
\end{proof}

\begin{example}\label{representation-theory:example-s3}
$S_3$ has three conjugacy classes: $\{e\}$, the three transpositions and the two $3$-cycles. Its irreducible
representations are the trivial representation, the sign representation and the two-dimensional
representation $U = \{x \in \CC^3 : x_1 + x_2 + x_3 = 0\}$ (a subrepresentation of the permutation representation).
Their characters take the values $(1, 1, 1)$, $(1, -1, 1)$ and $(2, 0, -1)$ on the three classes. One checks
orthonormality directly: for instance $\langle \chi_U, \chi_U\rangle = \frac16(4 + 0 + 2 \cdot 1) = 1$.
\end{example}

\begin{proposition}\label{representation-theory:proposition-abelian}
Let $G$ be a finite abelian group. Every irreducible complex representation of $G$ is one-dimensional, given by a homomorphism $\chi \colon G \to \CC^\times$, and there are exactly $|G|$ of them. For $G = \ZZ/n\ZZ$ they are
$\chi_j(k) = e^{2\pi ijk/n}$, $0 \leq j < n$.
\end{proposition}

\begin{proof}
Let $V$ be irreducible. Each $g$ commutes with the action of $G$, so by Schur's lemma (Lemma~\ref{representation-theory:lemma-schur}) it acts by a scalar; then every line in $V$ is a subrepresentation, and $\dim V = 1$. The conjugacy
classes of an abelian group are its elements, so by Theorem~\ref{representation-theory:theorem-number-irreducibles} there are $|G|$ irreducible representations. For $\ZZ/n\ZZ$, a homomorphism to $\CC^\times$ is determined by the image
$\zeta$ of $1$, with $\zeta^n = 1$, and the $n$ choices give the $\chi_j$.
\end{proof}

\begin{counterexample}\label{representation-theory:counterexample-schur-real}
The second statement of Schur's lemma needs an algebraically closed field. The group $\ZZ/4\ZZ$ acting on $\RR^2$ by rotation by $\pi/2$, with generator acting by $J = \begin{pmatrix}0 & -1\\ 1 & 0\end{pmatrix}$, is irreducible over $\RR$,
since $J$ has no real eigenvector; but its endomorphisms, the matrices commuting with $J$, are the $aI + bJ$, a field isomorphic to $\CC$, not $\RR\cdot\mathrm{id}$. Over $\CC$ the representation splits as $\chi_1 \oplus \chi_3$ in the notation of
Proposition~\ref{representation-theory:proposition-abelian}.
\end{counterexample}

\begin{definition}\label{representation-theory:definition-induced}
Let $H \subseteq G$ be a subgroup of a finite group and $W$ a representation of $H$. The \emph{induced representation} is
$\operatorname{Ind}_H^GW = \{f \colon G \to W : f(hx) = hf(x) \text{ for } h \in H, x \in G\}$, with $(gf)(x) = f(xg)$. For a representation $V$ of $G$, $\operatorname{Res}^G_HV$ is $V$ viewed as a representation of $H$.
\end{definition}

\begin{proposition}[Frobenius reciprocity]\label{representation-theory:proposition-frobenius}
In the situation of Definition~\ref{representation-theory:definition-induced}:
\begin{enumerate}
\item $\Hom_G(V, \operatorname{Ind}_H^GW) \cong \Hom_H(\operatorname{Res}^G_HV, W)$, via $\varphi \mapsto (v \mapsto \varphi(v)(e))$;
\item $\dim\operatorname{Ind}_H^GW = [G : H]\dim W$, and
\[
\chi_{\operatorname{Ind}W}(g) = \frac{1}{|H|}\sum_{x \in G,\ xgx^{-1} \in H}\chi_W(xgx^{-1});
\]
\item $\langle\chi_{\operatorname{Ind}W}, \chi_V\rangle_G = \langle\chi_W, \chi_{\operatorname{Res}V}\rangle_H$.
\end{enumerate}
\end{proposition}

\begin{proof}
(1) If $\varphi$ is $G$-equivariant, then $\psi(v) = \varphi(v)(e)$ satisfies $\psi(hv) = \varphi(hv)(e) = (h\varphi(v))(e) = \varphi(v)(h) = h\varphi(v)(e) = h\psi(v)$. Conversely, for $\psi \in \Hom_H(V, W)$ put $\varphi(v)(x) = \psi(xv)$; then
$\varphi(v)(hx) = \psi(hxv) = h\varphi(v)(x)$, so $\varphi(v) \in \operatorname{Ind}W$, and $\varphi(gv)(x) = \psi(xgv) = \varphi(v)(xg) = (g\varphi(v))(x)$. The two constructions are inverse, since $\varphi(v)(x) = (x\varphi(v))(e) = \varphi(xv)(e)$.

(2) Choose representatives $x_1, \ldots, x_m$ of the cosets $Hx$, $m = [G : H]$. An element $f$ is determined by the values $f(x_i)$, which are arbitrary, so $f \mapsto (f(x_i))_i$ is an isomorphism with $W^m$. For $g \in G$ write
$x_ig = h_ix_{\sigma(i)}$ with $h_i \in H$ and $\sigma$ a permutation; then $(gf)(x_i) = f(x_ig) = h_if(x_{\sigma(i)})$. So in the decomposition $W^m$, $g$ maps the $\sigma(i)$-th summand to the $i$-th by $h_i$, and only the indices with
$\sigma(i) = i$ contribute to the trace, with $h_i = x_igx_i^{-1} \in H$. Hence $\chi_{\operatorname{Ind}W}(g) = \sum_{i : x_igx_i^{-1} \in H}\chi_W(x_igx_i^{-1})$, and since $\chi_W$ is a class function on $H$, each coset $Hx_i$ contributes $|H|$ equal terms to
the sum in the statement.

(3) By Theorem~\ref{representation-theory:theorem-orthogonality}, both sides are dimensions of the spaces in (1): the left side is $\dim\Hom_G(V, \operatorname{Ind}_H^GW)$ and the right side is $\dim\Hom_H(\operatorname{Res}^G_HV, W)$.
\end{proof}

\begin{counterexample}\label{representation-theory:counterexample-character-table}
The character table does not determine the group. Let $G$ be a group of order $8$ whose centre $Z$ has order $2$ and with $G/Z \cong (\ZZ/2\ZZ)^2$; the dihedral group $D_4$ of symmetries of a square and the quaternion group
$Q_8 = \{\pm1, \pm i, \pm j, \pm k\}$ are examples. For $x \notin Z$ the centraliser contains $\langle x\rangle Z$, of order $4$, and is not $G$, so the class of $x$ is $\{x, zx\}$, $Z = \{e, z\}$: there are five classes. For $x \notin Z$, $x^2 \in Z$ and every conjugate of $x$ lies in $xZ$, because $G/Z$ is abelian of exponent $2$. The commutator subgroup is
$Z$ ($G/Z$ is abelian and $G$ is not), so the one-dimensional characters are the four characters of $G/Z$, equal to $1$ at $z$; at each non-identity element of $G/Z$ two of them take the value $1$ and two the value $-1$; by Theorem~\ref{representation-theory:theorem-number-irreducibles} the fifth irreducible character $\chi$ has
degree $2$ ($4 + 2^2 = 8$), and comparing with the regular character, $2\chi = \chi_{\mathrm{reg}} - \sum\chi_{1\text{-dim}}$ gives $\chi(z) = -2$ and $\chi(x) = 0$ for $x \notin Z$. So the character tables of $D_4$ and $Q_8$ coincide, but the
groups are not isomorphic: $D_4$ has five elements of order $2$ and $Q_8$ only one.
\end{counterexample}

\section{Lie algebras}\label{representation-theory:section-lie-algebras}

\begin{definition}\label{representation-theory:definition-lie-algebra}
A \emph{Lie algebra} over $k$ is a vector space $\mathfrak{g}$ with a bilinear map
$[-, -] \colon \mathfrak{g} \times \mathfrak{g} \to \mathfrak{g}$, the \emph{bracket}, with $[x, x] = 0$ and the
\emph{Jacobi identity} $[x, [y, z]] + [y, [z, x]] + [z, [x, y]] = 0$. A \emph{homomorphism} of Lie algebras is a
linear map preserving brackets, an \emph{ideal} is a subspace $\mathfrak{a}$ with
$[\mathfrak{g}, \mathfrak{a}] \subseteq \mathfrak{a}$, and $\mathfrak{g}$ is \emph{abelian} if its bracket is zero.
\end{definition}

\begin{example}\label{representation-theory:example-lie-algebras}
\begin{enumerate}
\item For an associative algebra $A$, the commutator $[x, y] = xy - yx$ makes $A$ a Lie algebra. For $A = \End(V)$
this is $\mathfrak{gl}(V)$, and for $A = M_n(k)$ it is $\mathfrak{gl}_n(k)$.
\item $\mathfrak{sl}_n(k) = \{x \in \mathfrak{gl}_n(k) : \operatorname{tr} x = 0\}$ is a Lie subalgebra, since
$\operatorname{tr}[x, y] = 0$ (Exercise~\ref{linear-algebra:exercise-trace}).
\item The skew-symmetric matrices form $\mathfrak{so}_n(k)$, and more generally the $x$ with $B(xv, w) + B(v, xw) = 0$
for a bilinear form $B$ form a Lie subalgebra of $\mathfrak{gl}(V)$.
\end{enumerate}
\end{example}

\begin{definition}\label{representation-theory:definition-lie-representation}
A \emph{representation} of a Lie algebra $\mathfrak{g}$ is a vector space $V$ with a Lie algebra homomorphism
$\mathfrak{g} \to \mathfrak{gl}(V)$; we write $xv$ for the action, so $[x, y]v = x(yv) - y(xv)$. Subrepresentations,
irreducibility, homomorphisms and direct sums are defined as for groups. The tensor product $V \otimes W$ is a
representation with $x(v \otimes w) = xv \otimes w + v \otimes xw$, and the dual with $(x\varphi)(v) = -\varphi(xv)$. The
\emph{adjoint representation} is $\mathfrak{g}$ itself with $x \cdot y = [x, y]$, a representation by the Jacobi
identity.
\end{definition}

\begin{remark}\label{representation-theory:remark-lie-groups-algebras}
Lie algebras arise as tangent spaces of Lie groups at the identity (Chapter~\ref{lie-groups}), and
representations of a Lie group differentiate to representations of its Lie algebra; the formula for the tensor
product is the derivative of $g(v \otimes w) = gv \otimes gw$.
\end{remark}

\section{Nilpotent and solvable Lie algebras}\label{representation-theory:section-solvable}

The Lie algebras of upper triangular and of strictly upper triangular matrices are the models of two classes of Lie algebras, the solvable and the nilpotent ones.
The theorems of Engel and Lie say that, conversely, a Lie algebra of nilpotent matrices can be put in strictly upper triangular form, and a solvable Lie algebra of
matrices over $\CC$ in upper triangular form, by one change of basis for all its elements at once. In this section all Lie algebras and representations are
finite-dimensional. For subspaces $\mathfrak{a}, \mathfrak{b} \subseteq \mathfrak{g}$, $[\mathfrak{a}, \mathfrak{b}]$ is the subspace spanned by the brackets $[x, y]$
with $x \in \mathfrak{a}$, $y \in \mathfrak{b}$, and the \emph{centre} of $\mathfrak{g}$ is $\mathfrak{z}(\mathfrak{g}) = \{x : [x, y] = 0 \text{ for all } y\}$. We write
$\ad x \colon \mathfrak{g} \to \mathfrak{g}$ for $y \mapsto [x, y]$, the adjoint representation of Definition~\ref{representation-theory:definition-lie-representation}.

\begin{definition}\label{representation-theory:definition-solvable}
The \emph{derived series} of a Lie algebra $\mathfrak{g}$ is $D^0\mathfrak{g} = \mathfrak{g}$, $D^{i+1}\mathfrak{g} = [D^i\mathfrak{g}, D^i\mathfrak{g}]$, and its
\emph{lower central series} is $C^1\mathfrak{g} = \mathfrak{g}$, $C^{i+1}\mathfrak{g} = [\mathfrak{g}, C^i\mathfrak{g}]$. The Lie algebra $\mathfrak{g}$ is \emph{solvable}
if $D^m\mathfrak{g} = 0$ for some $m$, and \emph{nilpotent} if $C^m\mathfrak{g} = 0$ for some $m$.
\end{definition}

So $\mathfrak{g}$ is nilpotent when all iterated brackets $[x_1, [x_2, [\cdots, [x_{m-1}, x_m]\cdots]]]$ of a fixed length $m$ vanish; abelian Lie algebras are nilpotent
with $m = 2$.

\begin{lemma}\label{representation-theory:lemma-solvable}
\begin{enumerate}
\item If $\mathfrak{a}$ and $\mathfrak{b}$ are ideals, so is $[\mathfrak{a}, \mathfrak{b}]$. Hence all $D^i\mathfrak{g}$ and $C^i\mathfrak{g}$ are ideals, and
$D^i\mathfrak{g} \subseteq C^{i+1}\mathfrak{g}$; in particular nilpotent Lie algebras are solvable.
\item Subalgebras and homomorphic images of solvable (nilpotent) Lie algebras are solvable (nilpotent).
\item If $\mathfrak{a}$ is a solvable ideal of $\mathfrak{g}$ and $\mathfrak{g}/\mathfrak{a}$ is solvable, then $\mathfrak{g}$ is solvable.
\item A nonzero nilpotent Lie algebra has nonzero centre; and if $\mathfrak{g}/\mathfrak{z}(\mathfrak{g})$ is nilpotent, so is $\mathfrak{g}$.
\end{enumerate}
\end{lemma}

\begin{proof}
(1) For $x \in \mathfrak{g}$, $a \in \mathfrak{a}$, $b \in \mathfrak{b}$ the Jacobi identity gives $[x, [a, b]] = [[x, a], b] + [a, [x, b]] \in [\mathfrak{a}, \mathfrak{b}]$.
By induction the terms of both series are ideals. Also $D^0\mathfrak{g} = C^1\mathfrak{g}$, and if $D^i\mathfrak{g} \subseteq C^{i+1}\mathfrak{g}$ then
$D^{i+1}\mathfrak{g} = [D^i\mathfrak{g}, D^i\mathfrak{g}] \subseteq [\mathfrak{g}, C^{i+1}\mathfrak{g}] = C^{i+2}\mathfrak{g}$.

(2) For a subalgebra $\mathfrak{h}$, induction gives $D^i\mathfrak{h} \subseteq D^i\mathfrak{g}$ and $C^i\mathfrak{h} \subseteq C^i\mathfrak{g}$. For a surjective
homomorphism $\varphi \colon \mathfrak{g} \to \mathfrak{g}^{\prime}$, $\varphi([\mathfrak{a}, \mathfrak{b}]) = [\varphi(\mathfrak{a}), \varphi(\mathfrak{b})]$, so
$\varphi(D^i\mathfrak{g}) = D^i\mathfrak{g}^{\prime}$ and $\varphi(C^i\mathfrak{g}) = C^i\mathfrak{g}^{\prime}$.

(3) If $D^m(\mathfrak{g}/\mathfrak{a}) = 0$, then $D^m\mathfrak{g} \subseteq \mathfrak{a}$ by (2) applied to the projection, and by induction
$D^{m+l}\mathfrak{g} \subseteq D^l\mathfrak{a}$, which vanishes for large $l$.

(4) Let $m$ be least with $C^m\mathfrak{g} = 0$; as $\mathfrak{g} \neq 0$, $m \geq 2$ and $C^{m-1}\mathfrak{g} \neq 0$. Since
$[\mathfrak{g}, C^{m-1}\mathfrak{g}] = C^m\mathfrak{g} = 0$, $C^{m-1}\mathfrak{g} \subseteq \mathfrak{z}(\mathfrak{g})$. If $C^m(\mathfrak{g}/\mathfrak{z}(\mathfrak{g})) = 0$,
then $C^m\mathfrak{g} \subseteq \mathfrak{z}(\mathfrak{g})$ and $C^{m+1}\mathfrak{g} = [\mathfrak{g}, C^m\mathfrak{g}] = 0$.
\end{proof}

\begin{example}\label{representation-theory:example-triangular}
Let $\mathfrak{b}_n \subseteq \mathfrak{gl}_n(k)$ be the upper triangular matrices, $\mathfrak{n}_n$ the strictly upper triangular ones and $\mathfrak{d}_n$ the diagonal
ones; all three are Lie subalgebras.
\begin{enumerate}
\item $\mathfrak{n}_n$ is nilpotent. Let $\mathfrak{n}^{(j)}$ be the matrices $(a_{il})$ with $a_{il} = 0$ unless $l - i \geq j$, so that
$\mathfrak{n}_n = \mathfrak{n}^{(1)}$ and $\mathfrak{n}^{(n)} = 0$. The entry $(AB)_{il} = \sum_ma_{im}b_{ml}$ of a product of $A \in \mathfrak{n}^{(j)}$ and
$B \in \mathfrak{n}^{(j^{\prime})}$ can be nonzero only if $m - i \geq j$ and $l - m \geq j^{\prime}$ for some $m$, so $AB \in \mathfrak{n}^{(j + j^{\prime})}$. Hence
$[\mathfrak{n}^{(1)}, \mathfrak{n}^{(j)}] \subseteq \mathfrak{n}^{(j+1)}$, and $C^i\mathfrak{n}_n \subseteq \mathfrak{n}^{(i)}$ vanishes for $i \geq n$.
\item $\mathfrak{b}_n$ is solvable: the diagonal of a product of upper triangular matrices is the product of the diagonals, so the diagonal of a commutator vanishes and
$[\mathfrak{b}_n, \mathfrak{b}_n] \subseteq \mathfrak{n}_n$, which is solvable by (1) and Lemma~\ref{representation-theory:lemma-solvable}(1); so
$D^{m+1}\mathfrak{b}_n \subseteq D^m\mathfrak{n}_n$ vanishes for large $m$. For $n \geq 2$ it is not nilpotent: with the matrix units $E_{il}$,
$[E_{11}, E_{12}] = E_{12}$, so $E_{12} \in C^i\mathfrak{b}_n$ for all $i$.
\item The two-dimensional Lie algebra with basis $x$, $y$ and $[x, y] = y$ (for instance $x = E_{11}$, $y = E_{12}$ in $\mathfrak{b}_2$) is solvable, with
$D^1 = ky$ and $D^2 = 0$, but not nilpotent, since $C^i = ky$ for all $i \geq 2$. The \emph{Heisenberg algebra} $\mathfrak{n}_3$, with basis $E_{12}, E_{23}, E_{13}$
and the only nonzero bracket $[E_{12}, E_{23}] = E_{13}$, is nilpotent but not abelian.
\item If $2 \neq 0$ in $k$, $\mathfrak{sl}_2(k)$ is not solvable: its basis $e = E_{12}$, $f = E_{21}$, $h = E_{11} - E_{22}$ satisfies
$[h, e] = 2e$, $[h, f] = -2f$ and $[e, f] = h$, so $[\mathfrak{sl}_2, \mathfrak{sl}_2] = \mathfrak{sl}_2$ and $D^i\mathfrak{sl}_2 = \mathfrak{sl}_2$ for all $i$.
\end{enumerate}
\end{example}

\begin{lemma}\label{representation-theory:lemma-ad-nilpotent}
If $x \in \mathfrak{gl}(V)$ is nilpotent, then $\ad x \colon \mathfrak{gl}(V) \to \mathfrak{gl}(V)$ is nilpotent.
\end{lemma}

\begin{proof}
Write $\ad x = \lambda_x - \rho_x$ with $\lambda_x(y) = xy$ and $\rho_x(y) = yx$. These two linear maps commute, by associativity, and if $x^m = 0$ then
$\lambda_x^m = \lambda_{x^m} = 0$ and $\rho_x^m = 0$. By the binomial formula for commuting maps, $(\lambda_x - \rho_x)^{2m} = \sum_i\binom{2m}{i}\lambda_x^i(-\rho_x)^{2m - i}$,
and every term has $i \geq m$ or $2m - i \geq m$, hence vanishes.
\end{proof}

\begin{theorem}[Engel's theorem, linear form]\label{representation-theory:theorem-engel-linear}
Let $V \neq 0$ be a finite-dimensional vector space over any field $k$, and $\mathfrak{g} \subseteq \mathfrak{gl}(V)$ a Lie subalgebra consisting of nilpotent
endomorphisms. Then there is $v \neq 0$ in $V$ with $xv = 0$ for all $x \in \mathfrak{g}$. Consequently $V$ has a basis in which every element of $\mathfrak{g}$ is
strictly upper triangular.
\end{theorem}

\begin{proof}
Induction on $\dim\mathfrak{g}$. If $\mathfrak{g} = 0$ any $v \neq 0$ works, and if $\dim\mathfrak{g} = 1$, $\mathfrak{g} = kx$ with $x$ nilpotent, and any nonzero
$v \in \ker x$ works ($\ker x \neq 0$, as $x$ is not invertible). Let $\dim\mathfrak{g} \geq 2$ and let $\mathfrak{h}$ be a proper subalgebra of maximal dimension
(for instance $0$ is proper).

\emph{Step 1: $\mathfrak{h}$ is an ideal of codimension $1$.} Since $[\mathfrak{h}, \mathfrak{h}] \subseteq \mathfrak{h}$, each $\ad h$, $h \in \mathfrak{h}$, preserves
$\mathfrak{h}$ and induces an endomorphism $\overline{\ad h}$ of the nonzero space $\mathfrak{g}/\mathfrak{h}$; by the Jacobi identity $h \mapsto \overline{\ad h}$ is a
homomorphism of Lie algebras $\mathfrak{h} \to \mathfrak{gl}(\mathfrak{g}/\mathfrak{h})$. Each $\ad h$ is nilpotent on $\mathfrak{gl}(V)$
(Lemma~\ref{representation-theory:lemma-ad-nilpotent}), hence on the invariant subspace $\mathfrak{g}$ and on the quotient $\mathfrak{g}/\mathfrak{h}$. So the image of
$\mathfrak{h}$ is a Lie algebra of nilpotent endomorphisms of $\mathfrak{g}/\mathfrak{h}$ of dimension at most $\dim\mathfrak{h} < \dim\mathfrak{g}$, and by induction
there is $y \in \mathfrak{g} \setminus \mathfrak{h}$ with $[h, y] \in \mathfrak{h}$ for all $h \in \mathfrak{h}$. Then $\mathfrak{h} + ky$ is a subalgebra strictly
containing $\mathfrak{h}$, so it is $\mathfrak{g}$ by maximality; hence $\mathfrak{g} = \mathfrak{h} \oplus ky$ and $\mathfrak{h}$ is an ideal.

\emph{Step 2: a common null vector.} By induction applied to $\mathfrak{h} \subseteq \mathfrak{gl}(V)$, the space $W = \{v \in V : hv = 0 \text{ for all } h \in \mathfrak{h}\}$
is nonzero. It is stable under $y$: for $w \in W$ and $h \in \mathfrak{h}$, $h(yw) = y(hw) + [h, y]w = 0$, as $[h, y] \in \mathfrak{h}$. The restriction of the nilpotent
$y$ to $W$ is nilpotent, so it kills some $v \neq 0$ in $W$, and then $xv = 0$ for all $x \in \mathfrak{h} + ky = \mathfrak{g}$.

\emph{Step 3: the basis.} Let $v_1$ be as above. Each $x \in \mathfrak{g}$ induces a nilpotent endomorphism of $V/kv_1$, and these form a Lie algebra of nilpotent
endomorphisms; by induction on $\dim V$, $V/kv_1$ has a basis $\bar v_2, \ldots, \bar v_d$ in which they are strictly upper triangular. Lifting it to $v_2, \ldots, v_d$,
we get $xv_1 = 0$ and $xv_j \in \mathrm{span}(v_1, \ldots, v_{j-1})$, which is the claim.
\end{proof}

\begin{corollary}[Engel's theorem]\label{representation-theory:corollary-engel}
A Lie algebra $\mathfrak{g}$ is nilpotent if and only if $\ad x$ is nilpotent for every $x \in \mathfrak{g}$.
\end{corollary}

\begin{proof}
If $C^m\mathfrak{g} = 0$, then $\ad x$ maps $C^i\mathfrak{g}$ into $C^{i+1}\mathfrak{g}$, so $(\ad x)^{m-1}\mathfrak{g} \subseteq C^m\mathfrak{g} = 0$. Conversely, the image
$\ad\mathfrak{g} \subseteq \mathfrak{gl}(\mathfrak{g})$ is a Lie algebra of nilpotent endomorphisms. Theorem~\ref{representation-theory:theorem-engel-linear} gives a basis
$v_1, \ldots, v_d$ of $\mathfrak{g}$ such that, with $V_j = \mathrm{span}(v_1, \ldots, v_j)$, $[\mathfrak{g}, V_j] \subseteq V_{j-1}$. By induction
$C^i\mathfrak{g} \subseteq V_{d - i + 1}$, since $C^1\mathfrak{g} = V_d$ and $C^{i+1}\mathfrak{g} = [\mathfrak{g}, C^i\mathfrak{g}] \subseteq [\mathfrak{g}, V_{d-i+1}] \subseteq V_{d-i}$.
So $C^{d+1}\mathfrak{g} = 0$.
\end{proof}

\begin{counterexample}\label{representation-theory:counterexample-engel}
\begin{enumerate}
\item A nilpotent Lie algebra of matrices need not consist of nilpotent matrices: the diagonal matrices $\mathfrak{d}_n$ form an abelian, hence nilpotent, Lie algebra
containing the identity. Engel's theorem is about the operators $\ad x$ on $\mathfrak{g}$, not about the matrices $x$ on $V$.
\item In Theorem~\ref{representation-theory:theorem-engel-linear} it matters that $\mathfrak{g}$ is closed under brackets. The matrices $e = E_{12}$ and $f = E_{21}$ of
$\mathfrak{sl}_2$ are nilpotent, but $\ker e = ke_1$ and $\ker f = ke_2$ meet only in $0$; the Lie algebra they generate contains $[e, f] = h$, which is not nilpotent.
\end{enumerate}
\end{counterexample}

For Lie's theorem we need eigenvectors, hence an algebraically closed field, and characteristic $0$. The key step is the following lemma.

\begin{lemma}[Invariance lemma]\label{representation-theory:lemma-invariance}
Let $k$ have characteristic $0$, let $\mathfrak{g} \subseteq \mathfrak{gl}(V)$ be a Lie subalgebra, $\mathfrak{a}$ an ideal of $\mathfrak{g}$ and
$\lambda \colon \mathfrak{a} \to k$ a linear form. Then $W = \{v \in V : av = \lambda(a)v \text{ for all } a \in \mathfrak{a}\}$ is stable under $\mathfrak{g}$.
\end{lemma}

\begin{proof}
Let $w \in W$ and $x \in \mathfrak{g}$. For $a \in \mathfrak{a}$, $a(xw) = x(aw) + [a, x]w = \lambda(a)xw + \lambda([a, x])w$, as $[a, x] \in \mathfrak{a}$. So it suffices
to show $\lambda([a, x]) = 0$ for all $a$, when $W \neq 0$. Fix $w \neq 0$ in $W$, let $W_i = \mathrm{span}(w, xw, \ldots, x^{i-1}w)$ (with $W_0 = 0$), and let $m \geq 1$ be
least with $W_{m+1} = W_m$; then $w, xw, \ldots, x^{m-1}w$ is a basis of $W_m$, and $xW_m \subseteq W_m$.

We claim that $ax^iw \in \lambda(a)x^iw + W_i$ for all $a \in \mathfrak{a}$ and $i \geq 0$. For $i = 0$ this is $aw = \lambda(a)w$. If it holds for $i - 1$ and all
$a$, then $ax^iw = x(ax^{i-1}w) + [a, x]x^{i-1}w$. By induction $ax^{i-1}w = \lambda(a)x^{i-1}w + u$ with $u \in W_{i-1}$, so
$x(ax^{i-1}w) \in \lambda(a)x^iw + W_i$; and $[a, x]x^{i-1}w \in \lambda([a, x])x^{i-1}w + W_{i-1} \subseteq W_i$. This proves the claim.

So each $a \in \mathfrak{a}$ maps $W_m$ into itself, by an upper triangular matrix in the basis $w, xw, \ldots, x^{m-1}w$ with all diagonal entries $\lambda(a)$, and
$\operatorname{tr}(a|_{W_m}) = m\lambda(a)$. Apply this to $[a, x] \in \mathfrak{a}$: since $a$ and $x$ both preserve $W_m$, $[a, x]|_{W_m}$ is the commutator of
$a|_{W_m}$ and $x|_{W_m}$, which has trace $0$. So $m\lambda([a, x]) = 0$, and $\lambda([a, x]) = 0$ as $m \neq 0$ in $k$.
\end{proof}

\begin{theorem}[Lie's theorem]\label{representation-theory:theorem-lie}
Let $k$ be algebraically closed of characteristic $0$, $V \neq 0$ a finite-dimensional vector space and $\mathfrak{g} \subseteq \mathfrak{gl}(V)$ a solvable Lie
subalgebra. Then the elements of $\mathfrak{g}$ have a common eigenvector: there are $v \neq 0$ and a linear form $\lambda$ on $\mathfrak{g}$ with $xv = \lambda(x)v$ for
all $x \in \mathfrak{g}$. Consequently $V$ has a basis in which every element of $\mathfrak{g}$ is upper triangular.
\end{theorem}

\begin{proof}
Induction on $\dim\mathfrak{g}$; for $\mathfrak{g} = 0$ take any $v \neq 0$ and $\lambda = 0$. Let $\mathfrak{g} \neq 0$. Then $[\mathfrak{g}, \mathfrak{g}] \neq \mathfrak{g}$,
since otherwise $D^i\mathfrak{g} = \mathfrak{g}$ for all $i$. Choose a subspace $\mathfrak{a}$ of codimension $1$ containing $[\mathfrak{g}, \mathfrak{g}]$. It is an ideal,
as $[\mathfrak{g}, \mathfrak{a}] \subseteq [\mathfrak{g}, \mathfrak{g}] \subseteq \mathfrak{a}$, and it is solvable (Lemma~\ref{representation-theory:lemma-solvable}(2)). By
induction there are $v_0 \neq 0$ and $\lambda \colon \mathfrak{a} \to k$ with $av_0 = \lambda(a)v_0$ for $a \in \mathfrak{a}$, so the space $W$ of
Lemma~\ref{representation-theory:lemma-invariance} is nonzero, and it is stable under $\mathfrak{g}$. Write $\mathfrak{g} = \mathfrak{a} \oplus kz$. As $k$ is algebraically
closed, $z|_W$ has an eigenvector $v \in W$, $zv = \mu v$. Then $v$ is an eigenvector of every element of $\mathfrak{g}$, with $\lambda$ extended by $\lambda(z) = \mu$.
The basis is obtained as in Step 3 of the proof of Theorem~\ref{representation-theory:theorem-engel-linear}, applying the first part to $V/kv$.
\end{proof}

\begin{corollary}\label{representation-theory:corollary-lie}
Let $k$ be algebraically closed of characteristic $0$ and $\mathfrak{g}$ solvable.
\begin{enumerate}
\item Every irreducible representation of $\mathfrak{g}$ is one-dimensional.
\item There is a chain of ideals $0 = \mathfrak{g}_0 \subset \mathfrak{g}_1 \subset \cdots \subset \mathfrak{g}_d = \mathfrak{g}$ with $\dim\mathfrak{g}_i = i$.
\item $[\mathfrak{g}, \mathfrak{g}]$ is nilpotent. Conversely, over any field, a Lie algebra $\mathfrak{g}$ with $[\mathfrak{g}, \mathfrak{g}]$ nilpotent is solvable.
\end{enumerate}
\end{corollary}

\begin{proof}
(1) The image of $\mathfrak{g}$ in $\mathfrak{gl}(V)$ is solvable (Lemma~\ref{representation-theory:lemma-solvable}(2)), so it has a common eigenvector $v$, and $kv$ is a
subrepresentation, hence equal to $V$. (2) Apply Theorem~\ref{representation-theory:theorem-lie} to $\ad\mathfrak{g} \subseteq \mathfrak{gl}(\mathfrak{g})$: in the
resulting basis $v_1, \ldots, v_d$, the spans $\mathfrak{g}_i = \mathrm{span}(v_1, \ldots, v_i)$ satisfy $[\mathfrak{g}, \mathfrak{g}_i] \subseteq \mathfrak{g}_i$. (3) In
the same basis every $\ad x$ is upper triangular, so every $\ad[x, y] = [\ad x, \ad y]$ is strictly upper triangular, as is every element of
$\ad[\mathfrak{g}, \mathfrak{g}]$; hence for $z \in [\mathfrak{g}, \mathfrak{g}]$, $\ad z$ is nilpotent on $\mathfrak{g}$, and so on the invariant subspace
$[\mathfrak{g}, \mathfrak{g}]$. By Corollary~\ref{representation-theory:corollary-engel}, $[\mathfrak{g}, \mathfrak{g}]$ is nilpotent. Conversely, if
$[\mathfrak{g}, \mathfrak{g}]$ is nilpotent it is solvable, and $\mathfrak{g}/[\mathfrak{g}, \mathfrak{g}]$ is abelian, so $\mathfrak{g}$ is solvable by
Lemma~\ref{representation-theory:lemma-solvable}(3).
\end{proof}

\begin{counterexample}\label{representation-theory:counterexample-lie}
Both hypotheses on $k$ are needed in Theorem~\ref{representation-theory:theorem-lie}.
\begin{enumerate}
\item Over $\RR$, the abelian Lie algebra spanned by the rotation generator $\bigl(\begin{smallmatrix} 0 & -1\\ 1 & 0\end{smallmatrix}\bigr)$ has no eigenvector in
$\RR^2$.
\item Let $k$ be algebraically closed of characteristic $p > 0$, and $V = k^p$ with basis $e_0, \ldots, e_{p-1}$, indices modulo $p$. Let $xe_i = ie_i$ and
$ye_i = e_{i+1}$. Then $[x, y]e_i = (i + 1)e_{i+1} - ie_{i+1} = ye_i$, so $[x, y] = y$ and $\mathfrak{g} = kx + ky$ is solvable
(Example~\ref{representation-theory:example-triangular}(3)). But $x$ and $y$ have no common eigenvector. An eigenvector $v = \sum_ic_ie_i$ of $y$ with eigenvalue $\mu$
satisfies $c_{i-1} = \mu c_i$ for all $i$, hence $c_i = \mu^{p}c_i$ and $\mu^p = 1$, that is $(\mu - 1)^p = 0$ and $\mu = 1$; so $v$ is a multiple of
$\sum_ie_i$, and $xv = \sum_iie_i$ is not a multiple of $v$ since $p \geq 2$. In the proof of Lemma~\ref{representation-theory:lemma-invariance} the argument
breaks down at $m\lambda([a, x]) = 0$, with $m = p$.
\end{enumerate}
\end{counterexample}

\begin{remark}\label{representation-theory:remark-cartan-criterion}
Every Lie algebra $\mathfrak{g}$ has a largest solvable ideal, the \emph{radical} (the sum of two solvable ideals is solvable by
Lemma~\ref{representation-theory:lemma-solvable}(3), as $(\mathfrak{a} + \mathfrak{b})/\mathfrak{a} \cong \mathfrak{b}/(\mathfrak{a} \cap \mathfrak{b})$), and
$\mathfrak{g}$ is \emph{semisimple} if its radical is $0$. Over a field of characteristic $0$, \emph{Cartan's criterion} characterises solvability by the trace form:
$\mathfrak{g} \subseteq \mathfrak{gl}(V)$ is solvable if and only if $\operatorname{tr}(xy) = 0$ for all $x \in \mathfrak{g}$ and $y \in [\mathfrak{g}, \mathfrak{g}]$;
and $\mathfrak{g}$ is semisimple if and only if its Killing form $\operatorname{tr}(\ad x\,\ad y)$ is nondegenerate. The proof of Cartan's criterion uses Lie's theorem
and the Jordan--Chevalley decomposition (Theorem~\ref{linear-algebra:theorem-jordan-chevalley}); see \cite[Sections 3--5]{humphreys-lie}. For example $\mathfrak{sl}_2(\CC)$ is
semisimple, and Theorem~\ref{representation-theory:theorem-sl2-semisimple} below is the first case of Weyl's theorem
(Remark~\ref{representation-theory:remark-semisimple-lie-algebras}).
\end{remark}

\section{Representations of $\mathfrak{sl}_2$}\label{representation-theory:section-sl2}

In this section $k = \CC$ and $\mathfrak{sl}_2 = \mathfrak{sl}_2(\CC)$, with basis
\[
e = \begin{pmatrix} 0 & 1 \\ 0 & 0 \end{pmatrix}, \qquad f = \begin{pmatrix} 0 & 0 \\ 1 & 0 \end{pmatrix}, \qquad
h = \begin{pmatrix} 1 & 0 \\ 0 & -1 \end{pmatrix},
\]
satisfying $[h, e] = 2e$, $[h, f] = -2f$ and $[e, f] = h$. All representations are finite-dimensional.

\begin{lemma}\label{representation-theory:lemma-sl2-weights}
Let $V$ be a representation and for $\lambda \in \CC$ let $V_\lambda = \ker(h - \lambda)$ and
$V^{(\lambda)} = \bigcup_N \ker(h - \lambda)^N$ (the generalised eigenspace). Then $e$ maps $V_\lambda$ to $V_{\lambda + 2}$ and
$V^{(\lambda)}$ to $V^{(\lambda + 2)}$, and $f$ maps them to $V_{\lambda - 2}$ and $V^{(\lambda - 2)}$. Moreover, for $k \geq 1$,
\[
e f^k = f^k e + k f^{k-1}(h - k + 1) \quad \text{as operators on } V.
\]
\end{lemma}

\begin{proof}
From $he = eh + 2e$ we get $(h - \lambda - 2)e = e(h - \lambda)$, hence $(h - \lambda - 2)^N e = e(h - \lambda)^N$; similarly for $f$.
The identity holds for $k = 1$ since $ef = fe + h$. If it holds for $k$, then using $hf = fh - 2f$,
$ef^{k+1} = f^k e f + kf^{k-1}(h - k + 1)f = f^k(fe + h) + kf^k(h - k - 1) = f^{k+1}e + (k+1)f^k(h - k)$.
\end{proof}

\begin{theorem}\label{representation-theory:theorem-sl2-irreducibles}
For every $m \geq 0$ there is an irreducible representation $V(m)$ of dimension $m + 1$ with a basis
$v_0, \ldots, v_m$ such that
\[
hv_j = (m - 2j)v_j, \qquad fv_j = (j + 1)v_{j+1}, \qquad ev_j = (m - j + 1)v_{j-1},
\]
with $v_{-1} = v_{m+1} = 0$. Every irreducible representation is isomorphic to exactly one $V(m)$. More precisely,
if $V$ is any representation and $v_0 \neq 0$ satisfies $ev_0 = 0$ and $hv_0 = \lambda v_0$, then $\lambda = m$ is a
non-negative integer and the vectors $v_j = f^jv_0/j!$ span a subrepresentation isomorphic to $V(m)$.
\end{theorem}

\begin{proof}
The formulas define a representation: it suffices to check the three bracket relations on each $v_j$. For
instance $(ef - fe)v_j = (j + 1)(m - j)v_j - j(m - j + 1)v_j = (m - 2j)v_j = hv_j$, and $[h, e] = 2e$, $[h, f] = -2f$
hold because $e$ raises and $f$ lowers the $h$-eigenvalue by $2$. It is irreducible: a nonzero subrepresentation
is $h$-stable, so it contains some $v_j$ (the eigenvalues $m - 2j$ are distinct), then $v_0$ by applying $e$
(the coefficients $m - j + 1$ are nonzero for $1 \leq j \leq m$), then all $v_j$ by applying $f$.

Now let $v_0$ be as in the last statement and $v_j = f^jv_0/j!$. By
Lemma~\ref{representation-theory:lemma-sl2-weights}, $hv_j = (\lambda - 2j)v_j$, and $fv_j = (j + 1)v_{j+1}$ by definition.
By the identity of the lemma, $ev_j = \frac{1}{j!}jf^{j-1}(h - j + 1)v_0 = (\lambda - j + 1)v_{j-1}$. The nonzero $v_j$ have
distinct $h$-eigenvalues, hence are linearly independent, so there is a largest $m$ with $v_m \neq 0$, and then
$v_{m+1} = 0$. Hence $0 = ev_{m+1} = (\lambda - m)v_m$, so $\lambda = m$, and the span of $v_0, \ldots, v_m$ is a
subrepresentation isomorphic to $V(m)$.

Every representation $V \neq 0$ contains such a $v_0$: $h$ has an eigenvector $w$, with eigenvalue $\mu$, and among
$w, ew, e^2w, \ldots$ (which lie in the distinct eigenspaces $V_{\mu + 2i}$ or vanish) there is a last nonzero one.
If $V$ is irreducible, the subrepresentation generated by $v_0$ is all of $V$, so $V \cong V(m)$ with
$m = \dim V - 1$.
\end{proof}

\begin{theorem}\label{representation-theory:theorem-sl2-semisimple}
Every finite-dimensional representation of $\mathfrak{sl}_2$ is completely reducible, and $h$ acts
diagonalisably with integer eigenvalues.
\end{theorem}

\begin{proof}
Let $C = ef + fe + \frac12h^2$ (the \emph{Casimir operator}). A direct computation with the relations shows that $C$
commutes with $e$, $f$ and $h$; for example
$[C, e] = e[f, e] + [f, e]e + \frac12(h[h, e] + [h, e]h) = -eh - he + eh + he = 0$. On $V(m)$, $C$ is scalar by
Lemma~\ref{representation-theory:lemma-schur}, and evaluating on $v_0$ gives
$C = m + \frac12m^2 = \frac12m(m + 2)$ (using $efv_0 = mv_0$ and $ev_0 = 0$). These scalars are distinct for different $m \geq 0$.

Since $C$ commutes with the action, $V$ is the direct sum of the generalised eigenspaces of $C$, each of which is
a subrepresentation (Lemma~\ref{linear-algebra:lemma-kernel}). By choosing an irreducible subrepresentation
and passing to quotients repeatedly, every representation has a chain of subrepresentations with irreducible
successive quotients, each isomorphic to some $V(m)$ by
Theorem~\ref{representation-theory:theorem-sl2-irreducibles}; $C$ acts on the quotient $V(m)$ by
$\frac12m(m+2)$. Hence on the generalised $c$-eigenspace, all these quotients are isomorphic to the same $V(m)$
(with $c = \frac12m(m+2)$). It therefore suffices to prove: if $V$ has such a chain
$0 = F_0 \subsetneq F_1 \subsetneq \cdots \subsetneq F_r = V$ with $F_i/F_{i-1} \cong V(m)$ for all $i$, then $V \cong V(m)^r$.

The eigenvalues of $h$ on $V$ are among $m, m - 2, \ldots, -m$, so $V^{(m + 2)} = 0$ and $e$ kills
$V^{(m)}$ (Lemma~\ref{representation-theory:lemma-sl2-weights}). Also $\dim V^{(m)} = r$, since each
quotient contributes one dimension. The map $f^m \colon V^{(m)} \to V^{(-m)}$ is injective: if $u \in V^{(m)}$ is nonzero, let
$i$ be least with $u \in F_i$; the image of $u$ in $F_i/F_{i-1} \cong V(m)$ is a nonzero vector of weight $m$, on which
$f^m$ is nonzero, so $f^mu \notin F_{i-1}$. Now $f^{m+1}u \in V^{(-m-2)} = 0$, and by
Lemma~\ref{representation-theory:lemma-sl2-weights},
$0 = ef^{m+1}u = f^{m+1}eu + (m + 1)f^m(h - m)u = (m + 1)f^m(h - m)u$. Since $(h - m)u \in V^{(m)}$ and $f^m$ is
injective there, $(h - m)u = 0$. So $h = m$ on $V^{(m)}$.

Let $u_1, \ldots, u_r$ be a basis of $V^{(m)} = V_m$. Each $u_i$ satisfies $eu_i = 0$ and $hu_i = mu_i$, so it generates a
subrepresentation $W_i \cong V(m)$ by Theorem~\ref{representation-theory:theorem-sl2-irreducibles}. The kernel $K$
of $\bigoplus_i W_i \to V$ is a subrepresentation whose weight-$m$ part is zero, since the $u_i$ are independent. If
$K \neq 0$, the last paragraph of the proof of Theorem~\ref{representation-theory:theorem-sl2-irreducibles},
applied to $K$, gives $v_0 \in K$ with $v_0 \neq 0$ and $ev_0 = 0$. But in $V(m)$ the kernel of $e$ is $\CC v_0$, since
$e(\sum_j c_jv_j) = \sum_{j \geq 1} c_j(m - j + 1)v_{j-1}$ vanishes only if $c_j = 0$ for $j \geq 1$; so in
$\bigoplus_i W_i \cong V(m)^r$ the kernel of $e$ is the weight-$m$ space, and $v_0$ would be a nonzero element of
$K$ of weight $m$. Hence $K = 0$, the sum is direct, and by dimension ($r(m + 1) = \dim V$) it is all of $V$. In each
$W_i$, $h$ is diagonal with integer eigenvalues.
\end{proof}

\begin{corollary}\label{representation-theory:corollary-lefschetz-structure}
Let $V$ be a finite-dimensional representation of $\mathfrak{sl}_2$, with weight spaces $V_\lambda = \ker(h - \lambda)$,
$\lambda \in \ZZ$. For every $j \geq 0$ the map $f^j \colon V_j \to V_{-j}$ is an isomorphism. Let
$P_j = V_j \cap \ker e$ (the \emph{primitive} vectors of weight $j$). Then
\[
V = \bigoplus_{j \geq 0}\ \bigoplus_{i = 0}^{j} f^i P_j,
\]
and $f^i$ is injective on $P_j$ for $0 \leq i \leq j$.
\end{corollary}

\begin{proof}
By Theorem~\ref{representation-theory:theorem-sl2-semisimple} it suffices to treat $V = V(m)$, where
$V_{m - 2i} = \CC v_i$, $f^j$ maps $v_{(m-j)/2}$ to a nonzero multiple of $v_{(m+j)/2}$ when $m - j$ is even and non-negative,
$P_m = \CC v_0$, $P_j = 0$ for $j \neq m$, and $f^iP_m = \CC v_i$.
\end{proof}

\begin{remark}\label{representation-theory:remark-semisimple-lie-algebras}
Theorem~\ref{representation-theory:theorem-sl2-semisimple} is a special case of Weyl's theorem: every
finite-dimensional representation of a semisimple Lie algebra over $\CC$ is completely reducible. The irreducible
representations are classified by their highest weights, and the semisimple Lie algebras themselves by Dynkin
diagrams. See \cite{humphreys-lie} and \cite{fulton-harris}.
\end{remark}

\section{Exercises}\label{representation-theory:section-exercises}

\begin{exercise}\label{representation-theory:exercise-symmetric-powers}
Let $\mathfrak{sl}_2$ act on $\CC^2$ by matrix multiplication. Show that $S^m(\CC^2) \cong V(m)$.
\end{exercise}

\begin{exercise}\label{representation-theory:exercise-clebsch-gordan}
Show that $V(m) \otimes V(n) \cong \bigoplus_{i=0}^{\min(m,n)} V(m + n - 2i)$. (Hint: count weights.)
\end{exercise}

\begin{solution}
The vectors $v_a \otimes w_b$ form a basis of weight vectors, of weight $(m - 2a) + (n - 2b)$. So the weight
$l = m + n - 2s$ has multiplicity $c(s) = |\{(a, b) : 0 \leq a \leq m, 0 \leq b \leq n, a + b = s\}|$, and for
$0 \leq s \leq (m+n)/2$ one finds $c(s) = \min(s, m, n) + 1$. By Theorem~\ref{representation-theory:theorem-sl2-semisimple}
the representation is a direct sum of $V(l)$'s, and the multiplicity of $V(l)$, $l \geq 0$, is
$\dim V_l - \dim V_{l+2} = c(s) - c(s - 1)$ (with $c(-1) = 0$), which is $1$ for $s \leq \min(m, n)$ and $0$ otherwise.
\end{solution}

\begin{exercise}\label{representation-theory:exercise-character-table-s4}
Determine the character table of $S_4$.
\end{exercise}

\begin{exercise}\label{representation-theory:exercise-abelian-irreducible}
Show that a finite group is abelian if and only if all its irreducible complex representations are one-dimensional. (One direction is
Proposition~\ref{representation-theory:proposition-abelian}; for the other use Theorem~\ref{representation-theory:theorem-number-irreducibles}.)
\end{exercise}

\begin{exercise}\label{representation-theory:exercise-sl2-lefschetz}
Let $V = \bigoplus_{k=0}^{2n} V^k$ be a graded finite-dimensional vector space and $L \colon V^k \to V^{k+2}$ a linear map. Suppose there is
$\Lambda \colon V^k \to V^{k-2}$ such that, with $H = (k - n)\id$ on $V^k$, the relations $[H, L] = 2L$, $[H, \Lambda] = -2\Lambda$,
$[L, \Lambda] = H$ hold. Show that $L^j \colon V^{n-j} \to V^{n+j}$ is an isomorphism for all $j \geq 0$. (This is the algebraic form
of the hard Lefschetz theorem.)
\end{exercise}

\begin{exercise}\label{representation-theory:exercise-two-dimensional}
Show that every Lie algebra of dimension $2$ is either abelian or isomorphic to the Lie algebra with basis $x$, $y$ and $[x, y] = y$ of
Example~\ref{representation-theory:example-triangular}(3).
\end{exercise}

\begin{solution}
Let $u$, $v$ be a basis. The bracket is determined by $[u, v]$, as $[u, u] = [v, v] = 0$. If $[u, v] = 0$ the algebra is abelian. Otherwise $y = [u, v]$ spans
$[\mathfrak{g}, \mathfrak{g}]$, which is one-dimensional. Complete $y$ to a basis $y$, $x^{\prime}$; then $[x^{\prime}, y] = cy$ for some $c \in k$, as $[x^{\prime}, y] \in [\mathfrak{g}, \mathfrak{g}]$,
and $c \neq 0$ because $[x^{\prime}, y]$ spans $[\mathfrak{g}, \mathfrak{g}]$. So $x = x^{\prime}/c$ satisfies $[x, y] = y$.
\end{solution}

\begin{exercise}\label{representation-theory:exercise-so3-not-solvable}
Let $\mathfrak{g}$ be the real Lie algebra $\RR^3$ with the cross product as bracket (it is isomorphic to the Lie algebra $\mathfrak{so}(3)$ of the rotation group,
Chapter~\ref{lie-groups}). Show that $[\mathfrak{g}, \mathfrak{g}] = \mathfrak{g}$, so that $\mathfrak{g}$ is not solvable, and that it has no ideals other than $0$ and itself.
\end{exercise}

\begin{solution}
$e_1 \times e_2 = e_3$, $e_2 \times e_3 = e_1$ and $e_3 \times e_1 = e_2$, so the brackets span everything and $D^i\mathfrak{g} = \mathfrak{g}$ for all $i$. The cross product satisfies the Jacobi identity, so $\mathfrak{g}$ is a Lie algebra. If $\mathfrak{a}$ is an
ideal containing some $a \neq 0$, then $\mathfrak{a} \supseteq a \times \RR^3$, and $a \times \RR^3 = a^\perp$, since $v \mapsto a \times v$ has kernel $\RR a$ and image in
$a^\perp$. For an orthonormal basis $u, w$ of $a^\perp$, $u \times w$ is a unit vector orthogonal to $u$ and $w$, so $u \times w = \pm a/|a|$ lies in $\mathfrak{a}$. Hence
$\mathfrak{a} \supseteq a^\perp + \RR a = \RR^3$.
\end{solution}

\begin{exercise}\label{representation-theory:exercise-killing-form-nilpotent}
Show that the Killing form $\operatorname{tr}(\ad x\,\ad y)$ of a nilpotent Lie algebra is zero.
\end{exercise}

\begin{solution}
By Corollary~\ref{representation-theory:corollary-engel} and Theorem~\ref{representation-theory:theorem-engel-linear} applied to $\ad\mathfrak{g}$, there is a basis of
$\mathfrak{g}$ in which every $\ad x$ is strictly upper triangular. A product of strictly upper triangular matrices is strictly upper triangular, so it has trace $0$.
\end{solution}
